\section*{Chapter \ref{ch:ml:from-prediction-to-portfolio} --- From Prediction to Portfolio}

\begin{solution}{exo:ml:from-prediction-to-portfolio:1}
Maximise $\mu w - \frac\gamma2\sigma^2w^2 - \kappa(w - w_{\mathrm{prev}})^2$: the first-order condition $\mu - \gamma\sigma^2w -
2\kappa(w - w_{\mathrm{prev}}) = 0$ gives the rule. With $\kappa = 0$ it is the mean-variance weight $\mu/(\gamma\sigma^2)$; as
$\kappa\to\infty$ the weight stays at $w_{\mathrm{prev}}$.
\end{solution}

\begin{solution}{exo:ml:from-prediction-to-portfolio:2}
The rule is $w = a\,\mu/(\gamma\sigma^2) + (1 - a)w_{\mathrm{prev}}$ with $a = \gamma\sigma^2/(\gamma\sigma^2 + 2\kappa) =
0.018/(0.018 + 0.001) = 0.947$: 95\% of the gap is closed each period; the cost proxy slows the portfolio only a little.
\end{solution}

\begin{solution}{exo:ml:from-prediction-to-portfolio:3}
Multiplying every return by a constant leaves the Sharpe ratio unchanged, so any scale of the coefficients with the same
direction is equally good up to the effect of costs; gradient ascent drifts in scale, larger positions trade more, and the
decision-focused turnover (2.75) exceeds aligned least squares' (1.59) at the same Sharpe ratio.
\end{solution}

\begin{solution}{exo:ml:from-prediction-to-portfolio:4}
Squared error is dominated by the untradeable names, where the first characteristic's signal is strong; fitting it lowers the
error there more than it raises it among the tradeable names, where that coefficient is pure error. Least squares on the
tradeable names fits only where it is scored by the portfolio.
\end{solution}

\begin{solution}{exo:ml:from-prediction-to-portfolio:5}
When the portfolio's sensitivity to the forecast differs across tradeable names or periods in a way a weight cannot express:
binding position limits for some names, costs that depend on trade size, a factor risk model that nets positions, a
forecast that feeds a non-linear rule. Then the errors that matter are not simply the errors on the traded names.
\end{solution}

\begin{solution}{exo:ml:from-prediction-to-portfolio:6}
The comparison is in the backtest the end-to-end model was fitted on, or selected on: a loss on the portfolio's Sharpe ratio
is a strong incentive to fit noise. Compare on held-out data, against a forecast trained on the traded names, with the spread
over independent samples; the chapter's small-sample case shows the end-to-end fit losing out of sample.
\end{solution}

\begin{solution}{exo:ml:from-prediction-to-portfolio:7}
Started from least squares on the tradeable names, with the number of \omterm{def:ml:neural-networks-for-noisy-tabular-data:backprop}{epochs} chosen on the last 60 training periods, the
validation block chooses no end-to-end training at all in three of the five samples; the mean test net Sharpe ratio is 1.04,
the same as aligned least squares. \omterm{def:ml:trees-and-boosting:early}{Early stopping} closes the gap by, mostly, stopping at the start.
\end{solution}

\begin{solution}{exo:ml:from-prediction-to-portfolio:8}
$w^* = \Sigma^{-1}\mu/\gamma$, so $\partial w^*/\partial\mu = \Sigma^{-1}/\gamma$. For a loss $L(w^*)$, $\partial L/\partial\mu =
(\partial w^*/\partial\mu)^\top\nabla_wL = \Sigma^{-1}\nabla_wL/\gamma$, and back-propagation carries it to the forecast's
parameters. The forecast's error matters in the directions the optimiser amplifies: those of low risk.
\end{solution}

\begin{solution}{pb:ml:from-prediction-to-portfolio:1}
\textbf{Part I.}
\begin{enumerate}
\item Untradeable names where the first characteristic predicts strongly, tradeable names where the second predicts weakly; a
single linear forecast cannot be right in both.
\item Mean-variance weights with a quadratic cost proxy on the tradeable names; net returns pay 5 basis points on every change
in weight.
\item 1.51 on the first characteristic and 0.43 on the second (0.1\% units): mostly the untradeable signal, which dominates the
squared error.
\item All names: 9.007 overall and 8.987 on the tradeable names; tradeable-names fit: 9.031 and 8.962 ($10^{-4}$ units).
\end{enumerate}
\textbf{Part II.}
\begin{enumerate}[resume]
\item Net, gross, turnover: 0.19, 0.60, 2.95; 1.52, 1.92, 1.59; 1.51, 1.92, 2.75; 1.49, 1.92, 2.90; truth 1.57, 1.97, 1.50.
\item 1.33 with a standard deviation of 0.18 over five pairs.
\item Both learn the second characteristic's direction; here the only thing the decision adds is which names matter, and a
restricted fit says the same.
\item It has no smoothing: weights move with the characteristics every period.
\end{enumerate}
\textbf{Part III.}
\begin{enumerate}[resume]
\item Aligned least squares 1.04, decision-focused 0.89, parametric 0.87, \omterm{def:ml:from-prediction-to-portfolio:pto}{predict-then-optimise} 0.09, truth 1.51.
\item The coefficients on useless characteristics, and turnover of 3.49 against 1.88.
\item Few parameters, a least-squares start, \omterm{def:ml:trees-and-boosting:early}{early stopping} on a validation block, penalties on the direction, not only the
scale.
\item On held-out periods and independent samples, by net performance and its spread, with aligned least squares as the
benchmark.
\end{enumerate}
\textbf{Part IV.}
\begin{enumerate}[resume]
\item \emph{The loss that pays.} \omterm{def:ml:from-prediction-to-portfolio:pto}{Predict-then-optimise} earns a net Sharpe ratio of 0.19, \omterm{def:ml:from-prediction-to-portfolio:dfl}{decision-focused learning} 1.51 and the
parametric \omterm{def:ml:reinforcement-learning-foundations:mdp}{policy} 1.49; the gap of 1.33 has a standard deviation of 0.18 over five train-test pairs, and least squares on the
traded names alone matches the end-to-end methods (1.52).
\item When the forecast model is well specified, or when it is fitted on the names and periods the decision uses.
\item Decisions whose sensitivity to the forecast varies in ways a sample restriction or weight cannot capture, and enough data
to fit through them.
\item The transfer coefficient measures what constraints lose between forecast and portfolio; this chapter's loss happens
before, when the forecast learns the wrong thing.
\item Its net portfolio performance under the desk's construction and costs, on held-out data, against the incumbent, with
turnover and spread.
\item In the portfolio's evaluation always; in training through a cost-aware rule or a cost-weighted loss.
\item The same, with the network in place of the linear forecast, and more regularisation.
\item To make the portfolio built from it good, net of costs, out of sample.
\end{enumerate}
\end{solution}

\begin{solution}{iq:ml:from-prediction-to-portfolio:1}
Accuracy is measured everywhere; profit depends on the forecast only where and as much as the optimiser trades on it, after
costs and constraints. A misspecified model trades errors between names, and the loss decides which.
\iqlookfor{misspecification and the mismatch between the loss and the decision.}
\end{solution}

\begin{solution}{iq:ml:from-prediction-to-portfolio:2}
Train the forecast by the loss of the decision made with it; differentiate the optimiser's solution in closed form, by
unrolling an iterative solver, or by implicit differentiation of its optimality conditions.
\iqlookfor{the loss and at least one differentiation method.}
\end{solution}

\begin{solution}{iq:ml:from-prediction-to-portfolio:3}
Weights are a linear function of characteristics around a benchmark, with coefficients chosen to maximise average utility in
sample. Few parameters, no covariance matrix, costs and constraints can enter the utility; but in-sample utility maximisation
overfits and the \omterm{def:ml:reinforcement-learning-foundations:mdp}{policy} form is restrictive.
\iqlookfor{the form, the fitting criterion, and the \omterm{def:ml:why-financial-machine-learning-is-different:generalisation}{overfitting} risk.}
\end{solution}

\begin{solution}{iq:ml:from-prediction-to-portfolio:4}
Evaluate the model by net portfolio returns; fit it on tradeable names, weighted by liquidity or expected position; or train
through a cost-aware portfolio rule; and penalise turnover in the forecast (smoothness) where costs are high.
\iqlookfor{net evaluation and at least one way to bring costs into training.}
\end{solution}

\begin{solution}{iq:ml:from-prediction-to-portfolio:5}
Whether it was evaluated on data used in fitting or selection, its turnover and costs, the spread over seeds and samples, the
benchmark (a forecast fitted on the traded names), and the number of parameters against the data.
\iqlookfor{held-out evaluation, a fair benchmark and a spread.}
\end{solution}

\begin{solution}{iq:ml:from-prediction-to-portfolio:6}
For $\min\frac12w^\top Qw - \mu^\top w$ subject to $Aw = b$, the KKT system $Qw - \mu + A^\top\lambda = 0$, $Aw = b$ is linear;
differentiating gives $\begin{pmatrix}Q & A^\top\\ A & 0\end{pmatrix}\begin{pmatrix}dw\\ d\lambda\end{pmatrix} =
\begin{pmatrix}d\mu\\ 0\end{pmatrix}$, so $dw/d\mu$ is the top-left block of the inverse; inequality constraints add the active
ones to $A$ (OptNet).
\iqlookfor{the linearised KKT system and active constraints.}
\end{solution}
